% Include-only section; requires amsmath, amssymb.
\section{Consecutive successful codes for an abstract biased parser}
\label{sec:biased-parser-synchronization}

This section proves an arithmetic completeness tool for a possible future
encoding.  It is not a description of the archived scratch-biased table,
whose source continuation is incorrect, and it is not needed to calibrate
ST215, whose candidate counter advances normally by one.  In particular,
no additional physical state bound or implementation claim follows here.

Let $P$ be Cantor pairing and $f$ its right projection.  Consider an
abstract parser with a bit $b\in\{0,1\}$, set to zero after S and to one
after MP or AX.  At a positive continuation pointer $L$, it decodes the
first field from $L-1+b$ and reads all further command fields by ordinary
right projection.  A fresh candidate is decoded without bias.  Use
\[
\begin{split}
 \operatorname{fields}(\mathsf S)&=[0],\\
 \operatorname{fields}(\mathsf{MP}(x))&=[1,x],\\
 \operatorname{fields}(\mathsf{AX}(c;x,y))&=[c+2,y,x].
\end{split}
\]
The retained arithmetic store may be arbitrary and need not be determined
by the pointer or bias.  Assume fresh-candidate numbers increase by either
one or two; the choice of increment may depend on the entire prior run.

\paragraph{A consecutive synchronization pair.}
For any $r>0$, put
\[
 a=P(2,P(1,P(0,r))),\qquad z=P(1,a),\qquad n=P(1,z).
\]
Fresh candidate $n$ parses
\[
                  \mathsf{MP}(1);\mathsf{AX}(0;0,1)
\]
and reaches pointer $r$ with bias one.  Indeed its first command leaves
pointer $a$ and sets bias one, so the next selector is decoded from $a$,
without an effective decrement.  The three AX fields are $[2,1,0]$.

For every $z\ge0$, the elementary Cantor identity is
\[
                         P(1,z)+1=P(0,z+1).
\]
Consequently fresh candidate $n+1$ begins with S, leaving pointer $z+1$
and bias zero.  Unless that S already halts, the next selector is read
from $z=P(1,a)$, so it is MP with argument two and remaining pointer
$P(1,P(0,r))$.  Its bias is one, and the next MP therefore has argument
zero and leaves pointer $r$.  Thus the second word is
\[
                       \mathsf S;\mathsf{MP}(2);\mathsf{MP}(0),
\]
and it too reaches pointer $r$, bias one, or halts sooner.

This synchronizes only the parser position and bias.  It does \emph{not}
assert equality of the two retained arithmetic stores.  A common suffix
must therefore begin with a reset before using a compiler whose input
requires zero counter and cache.

\paragraph{Positive-tail word encoding.}
Let $F(\vec a,r)$ be right-nested pairing of a finite field list over
residual $r$, and let $\beta(k)$ be the bias produced by command $k$.
For a finite command word $V$ and a positive unused tail $q$, define
\[
\begin{split}
 E_b([],q)&=q,\\
 E_b(k::V,q)&=1-b+
       F(\operatorname{fields}(k),E_{\beta(k)}(V,q)).
\end{split}
\]
Pairing dominates its right argument, so induction gives
$E_b(V,q)\ge q>0$.  From this pointer and bias $b$, the parser subtracts
$1-b$ and obtains exactly the nested fields of $k$.  It then reaches
the suffix pointer with precisely the bias required by the induction
hypothesis.  Thus it executes all of $V$, unless a command halts earlier,
and leaves $q$ if the whole word returns.

The positivity of the unused tail is necessary.  With base tail zero,
one would obtain $E_1([\mathsf S],0)=0$, which the scheduler would treat
as a fresh fetch; the intended final S would never execute.  The positive
tail avoids this zero alias without changing any command field.

\paragraph{Cofinal successful pairs.}
Suppose a strong refutation block $B_1;H$ halts from every
$(t,w,0,0)$, as in the inherited short-transfer compiler.  Let $V$ be
an AX reset followed by that block.  It is successful from every retained
store.  For each $q\ge1$, take $r=E_1(V,q)$ and construct the consecutive
pair $n(r),n(r)+1$ above.  The first synchronization word reaches $V$
from some retained store; the second either halts or does the same.
Therefore both codes are successful from every initial retained store.
Moreover
\[
                           n(r)\ge r\ge q,
\]
so these successful consecutive pairs are cofinal.

Every strictly increasing integer sequence with gaps at most two meets
every sufficiently late consecutive pair.  Otherwise, on first crossing
$n$, it would jump from a value at most $n-1$ to a value at least $n+2$,
a gap of at least three.  This argument does not assume that the gap
choices are independent of the candidate, parser bias or arithmetic data.
Thus skipping at most one candidate cannot defeat the successful pairs.

Each nonhalting candidate also drains.  For $L>0$, the first unpairing
input $v=L-1+b$ satisfies $0\le v\le L$.  If $v=0$, its right projection
is zero and is less than $L$.  If $v>0$, then $f(v)<v\le L$; equality
$f(v)=v$ would force the triangular part of $P$ to vanish and hence
$v=0$.  Further command fields only decrease the residual.  Every step
therefore strictly decreases the positive continuation pointer.  The
scheduler reaches each subsequent selected candidate in finitely many
commands unless it has already halted.

These are all-natural arithmetic arguments.  The companion checker
\texttt{check\_biased\_sync.py} corroborates the two
parsed prefixes for 1,000 positive residuals and checks 18 positive-tail
word encodings.  Those finite checks do not establish an implementation,
protected-tail soundness, or a physical state count.  An implementation
using this lemma would still need those separate obligations.
